% Geométrie de la page
\documentclass[a4paper,12pt,twoside]{article}
% \documentclass[a4paper,12pt]{article}
\usepackage[T1]{fontenc}
\usepackage[a4paper,top=2.5cm,bottom=3cm]{geometry}
\usepackage{fancyhdr}
\usepackage{listings}
\setlength{\headheight}{14.5pt}
\usepackage[french]{babel}
\usepackage{lmodern}
% \usepackage{graphicx} % Required for inserting images
\usepackage{tikz}
\usetikzlibrary{shapes.geometric, arrows.meta, positioning}

% Styles pour les schémas de flux
\tikzset{
    startstop/.style={rectangle, rounded corners, minimum width=2.8cm, minimum height=1cm, text centered, draw=black, fill=red!20},
    decision/.style={diamond, aspect=2, minimum width=3cm, minimum height=1cm, text centered, draw=black, fill=green!20},
    process/.style={rectangle, minimum width=2.8cm, minimum height=1cm, text centered, draw=black, fill=blue!20},
    fleche/.style={thick,->,>=Stealth}
}

% Propriétés du document
\author{Jonathan Vasseur}
\title{Résumé python 1M}

% % Headers and feet
\pagestyle{fancy}
\fancyhead{} % clear all header fields
\fancyhead[CO,CE]{\textbf{1M-\the\year{}}}
\fancyfoot{} % clear all footer fields
\fancyfoot[LE,RO]{\thepage}
% Groupes et commande custom
\newenvironment{ltitle}
     {\begin{center}\Huge}
     {\end{center}\par}

% % Début du document.
\begin{document}
    \begin{ltitle}
Résumé Python 1M
    \end{ltitle}

% % % % % % % % % % % % % % % % % % % % % % % % % % % % % 
% %                        Chapitre 1                 % % 
% % % % % % % % % % % % % % % % % % % % % % % % % % % % % 
    \section*{Types}
    \begin{center}
    
    \begin{tabular}{||c|c|c|c|c||}
    % \multicolumn{5}{||c||}{\Large{Types}} \\
    \hline \hline
                &\textbf{Integer}  & \textbf{Float} & \textbf{String} & \textbf{Boolean} \\
    \hline
    Diminutif   & Int   & Float  & Str  & Bool \\
    \hline
Exemple     & 1     & 1.0    & "Bonjour" & True \\
            & 0     & -2.5   & ""     & False   \\
            & -5    & .5     & "123"  &         \\
            & 45678 & 2.     & "123.5"&         \\
            &       &1e5     &        &         \\
    \hline \hline
    \end{tabular}
    \end{center}
% % % % % % % % % % % % % % % % % % % % % % % % % % % % % 
% %                        Chapitre 2                 % % 
% % % % % % % % % % % % % % % % % % % % % % % % % % % % % 
    \section*{Opérateurs et opérations}
    \subsection*{Mathématiques}
    \begin{tabular}{l c | l c}
    Addition: & a+b &
    Soustraction: & a-b \\ 
    Division: & a/b &
    Multiplication: & a*b  \\
    Division entière: & a//b &
    Modulo: & a\%b \\
    Puissance:& a**b 
    \end{tabular}
    \subsection*{Logiques}
    \begin{tabular}{l c | l c}
    Plus grand que: & a > b &
    Plus petit que: & a < b \\
    Plus grand ou égal que: & a >= b &
    Plus petit ou égal que: & a <= b \\
    Égale à: & a == b &
    Différent de: & a != b \\
    Inverse: & not a &
    Et logique: & a and b \\
    Ou logique: & a or b &  \\
    \end{tabular}
    \subsection*{Affectation composée}
    \begin{tabular}{l c | l c}
    Addition: & a += b &
    Soustraction: & a -= b \\
    Multiplication: & a *= b &
    Division: & a /= b \\
    Division entière: & a //= b &
    Modulo: & a \%= b \\
    Puissance:& a **= b
    \end{tabular}
    \subsection*{Commentaires}
    Une ligne précédée de \texttt{\#} est ignorée par Python, elle sert à documenter le code.
    \begin{lstlisting}[language=Python]
    # Ceci est un commentaire
    \end{lstlisting}
    \newpage

% % % % % % % % % % % % % % % % % % % % % % % % % % % % % 
% %                        Chapitre 3                 % % 
% % % % % % % % % % % % % % % % % % % % % % % % % % % % % 
\section*{Entrée}
    \subsection*{Fonction input}
    \begin{lstlisting}
    ma_chaine = input("Entrer du texte:\n")
    mon_nombre_entier = int(input("Entrer un nombre entier:\n"))
    mon_nombre_reel = float(input("Entrer un nombre reel:\n"))
    \end{lstlisting}

% % % % % % % % % % % % % % % % % % % % % % % % % % % % % 
% %                        Chapitre 4                 % % 
% % % % % % % % % % % % % % % % % % % % % % % % % % % % % 
    \section*{Sorties}
        \subsection*{Affichage et formatage des chaines de caractères}
        \begin{lstlisting}[language=Python]
    print(valeur1, valeur2, ..., sep=' ', end='\n')
        \end{lstlisting}
        Le paramètre \texttt{sep} est la chaîne de caractères insérée entre chaque
        valeur, par défaut un espace.\\
        Le paramètre \texttt{end} est la chaîne de caractères insérée après la
        dernière valeur, par défaut un retour de ligne.\\

        \subsection*{Technique de concaténation}
        \begin{lstlisting}[language=Python]
    chaine_concatenee = "Du texte" + chaine1 + str(nbre1) + chaine2 + ...
        \end{lstlisting} 

        \subsection*{F-string}
        \begin{lstlisting}[language=Python]
    f"blablabla {chaine1} blablabla {nbre1}"
        \end{lstlisting} 
    \newpage
% % % % % % % % % % % % % % % % % % % % % % % % % % % % % 
% %                        Chapitre 5                 % % 
% % % % % % % % % % % % % % % % % % % % % % % % % % % % % 
    \section*{Expressions conditionnelles}
        \subsection*{Syntaxe}
        \noindent Un bloc indenté n'est executé que \textbf{si et seulement si} la condition dont il dépend est \textbf{verifiée}.\\
        \begin{lstlisting}[language=Python]
    if <condition 1>:
        Instruction
        Instruction
        ...
    elif <condition 2>:
        Instruction
        Instruction
        ...
    else:
        Instruction
        Instruction
        ...
        \end{lstlisting} 
    \subsection*{Diagramme de flux}
    \begin{center}
    \begin{tikzpicture}[node distance=1.3cm and 2.2cm, every node/.style={font=\small}]
        \node (start) [startstop] {Début};
        \node (cond1) [decision, below=of start] {condition 1 ?};
        \node (instr1) [process, right=of cond1] {Instruction(s)};
        \node (cond2) [decision, below=of cond1] {condition 2 ?};
        \node (instr2) [process, right=of cond2] {Instruction(s)};
        \node (instr3) [process, below=of cond2] {Instruction(s) (else)};
        \node (end) [startstop, below=of instr3] {Fin};

        \draw [fleche] (start) -- (cond1);
        \draw [fleche] (cond1) -- node[above] {Vrai} (instr1);
        \draw [fleche] (cond1) -- node[left] {Faux} (cond2);
        \draw [fleche] (cond2) -- node[above] {Vrai} (instr2);
        \draw [fleche] (cond2) -- node[left] {Faux} (instr3);
        % contournement à droite pour ne pas traverser la boite instr2
        \draw [fleche] (instr1.east) -- ++(0.8,0) |- (end.east);
        \draw [fleche] (instr2.east) -- ++(0.4,0) |- (end.east);
        \draw [fleche] (instr3) -- (end);
    \end{tikzpicture}
    \end{center}
\newpage
% % % % % % % % % % % % % % % % % % % % % % % % % % % % % 
% %                        Chapitre 6                 % % 
% % % % % % % % % % % % % % % % % % % % % % % % % % % % % 
    \section*{Les boucles}
    \subsection*{La fonction range():}
    La fonction \texttt{range()} permet de crée une séquence d'entier selon le modèle suivant:
    \begin{lstlisting}[language=Python]
    range(start*,stop,step*)
    \end{lstlisting} 
    \texttt{start}: Valeur initiale de la séquence (par défaut 0)
    \texttt{stop}: Valeur limite d'arrêt (exclue)
    \texttt{step}: Valeur du pas de la séquence (par défaut 1)
    * paramètre optionnel

    \subsection*{for}
    \begin{lstlisting}[language=Python]
        for <variable> in <sequence>:
            Instruction
            Instruction
        \end{lstlisting} 
    
    \subsection*{while}
    \begin{lstlisting}[language=Python]
    <Initialisation de la condition>
    while <condition>:
        Instruction
        Instruction
        ...
    \end{lstlisting} 
    Quelque part à l'intérieur du \texttt{while} la condition doit évoluer.
    \end{document}


